\section{An integrated curvature certificate for the optimal vertex}
\label{sec:vertex}
The current FW gap contains more information than one oriented Bregman
divergence. Retaining the curvature profile along a segment to a
minimizer gives an observable gradient ball. The constant-profile
mechanism is the established cocoercivity and safe-screening
argument~\cite{wachsmuth,raj,sunbach};
the refinement here is the explicit integrated loss profile under
local-domain assumptions.

Let $f$ be convex and $C^2$ on an open convex neighborhood of compact
convex $Q$, and suppose
\begin{equation}\label{vertex:profile}
 0\preceq\nabla^2 f(z)\preceq\rho(\delta(z))I\qquad(z\in Q),
\end{equation}
where $\rho:[0,\infty)\to[0,\infty)$ is continuous and
nondecreasing. Define
\begin{equation}\label{vertex:profiledefs}
 \Omega(u)=\int_0^u\rho(v)\,dv,\qquad
 A(u)=\begin{cases}\Omega(u)/u,&u>0,\\ \rho(0),&u=0.\end{cases}
\end{equation}
The function $A$ is nondecreasing. The gradients at all constrained
minimizers coincide: along a segment joining two minimizers, convexity
forces the function to be constant; positive semidefiniteness of the
Hessian then gives $\nabla^2 f(z)d=0$ along the segment direction $d$,
so integration gives equal endpoint gradients. Denote this common value
by $g^*$.

\begin{proposition}[Integrated gradient-gap radius]
\label{vertex:radius}
For all $x\in Q$,
\begin{equation}\label{vertex:gradientbound}
 \norm{\nabla f(x)-g^*}^2\le A(\delta(x))G(x).
\end{equation}
In particular, valid bounds $\delta(x)\le U$ and $G(x)\le B$ give
$\norm{\nabla f(x)-g^*}\le\sqrt{A(U)B}$, and
$\norm{\nabla f(x)-g^*}^2\le\Omega(G(x))$. For
$\rho(u)=H_0+H_1u$, the coefficient is
$A(U)=H_0+H_1U/2$.
\end{proposition}
\begin{proof}
Fix a minimizer $y$ and put $d=x-y$ and
$H_t=\nabla^2 f(y+td)$ for $0\le t\le1$. Convexity gives
$0\le\delta(y+td)\le t\delta(x)$, whence
$0\preceq H_t\preceq\rho(t\delta(x))I$ by monotonicity of $\rho$.
The spectral inequality
$H_t^2\preceq\rho(t\delta(x))H_t$ and Cauchy--Schwarz imply
\begin{align*}
 \norm{\nabla f(x)-g^*}
 &\le\int_0^1\sqrt{\rho(t\delta(x))}
             \sqrt{d^TH_td}\,dt\\
 &\le\sqrt{A(\delta(x))\inner{\nabla f(x)-g^*}{d}}.
\end{align*}
Optimality gives $\inner{g^*}{d}\ge0$, and the feasible point $y$
in the FW gap gives $\inner{\nabla f(x)}{d}\le G(x)$. This proves
\eqref{vertex:gradientbound}. At $\delta(x)=0$ the gradient difference
is zero by the common-gradient property. Monotonicity of $A$ proves
the first consequence. Since convexity gives $\delta(x)\le G(x)$,
\eqref{vertex:gradientbound} also gives
$A(\delta(x))G(x)\le A(G(x))G(x)=\Omega(G(x))$.
\end{proof}

For a constant profile $\rho=L$, the coefficient one in
$\norm{\nabla f-g^*}^2\le LG$ is exact, as shown by
$f(t)=Lt^2/2$ on $[0,D]$. We do not claim that $\Omega(G)$ is the
sharp envelope for every nonconstant profile. For $\rho(u)=Hu^p$,
the integral gives $A(U)=HU^p/(p+1)$; compared with using $\rho(U)B$,
this reduces that sufficient squared radius by a factor $p+1$.

\begin{proposition}[A certified runner-up test]
\label{vertex:certificate}
Let $Q=\operatorname{conv}(V)$ have distinct finitely many vertices,
and suppose $\norm{\widehat g-\nabla f(x)}\le r$. Suppose
$s\in V$ satisfies the additive linear-oracle specification
\begin{equation}\label{vertex:oracle}
 \inner{\widehat g}{s}\le\min_{v\in V}\inner{\widehat g}{v}+\eta.
\end{equation}
Put
\begin{equation}\label{vertex:oraclebound}
 h=\inner{\widehat g}{x-s},\qquad B=h+\eta+Dr,
\end{equation}
and choose any valid $0\le U\le B$ with $\delta(x)\le U$. If no
tighter primal bound is available, use $U=B$. Set
$R=r+\sqrt{A(U)B}$. If a certified lower bound $\kappa$ satisfies
\begin{equation}\label{vertex:margin}
 \inner{\widehat g}{v-s}\ge\kappa\quad(v\in V\setminus\{s\}),
 \qquad \kappa>DR,
\end{equation}
then $X^*=\{s\}$.
\end{proposition}
\begin{proof}
The oracle specification gives
$\inner{\widehat g}{x-s}\ge-\eta$, since $x\in Q$ is a convex
combination of vertices. Moreover,
\[
 G(x)=\max_{v\in V}\inner{\nabla f(x)}{x-v}
 \le \inner{\widehat g}{x-s}+\eta+Dr
 =h+\eta+Dr=B.
\]
Thus Proposition~\ref{vertex:radius} and the gradient-error bound
give $\norm{\widehat g-g^*}\le R$. For every other vertex,
\[
 \inner{g^*}{v-s}
 \ge\inner{\widehat g}{v-s}-D\norm{g^*-\widehat g}
 \ge\kappa-DR>0.
\]
The linear form defined by $g^*$ therefore has the unique minimizer
$s$ on $Q$. Every minimizer of $f$ minimizes this linear form, so
$X^*=\{s\}$.
\end{proof}

A standard LMO alone does not provide $\kappa$. Computing it requires
an exhaustive vertex scan or a certified runner-up linear oracle; its
cost must be counted. A gap to merely sampled alternatives is
insufficient. If exact oracles are used, $G_k\to0$, and the problem has
a strictly exposed unique optimal vertex, the exact runner-up gap
converges to its positive optimal margin while $R_k\to0$. The test
eventually succeeds. An arbitrary conservative lower bound, such as
$\kappa\equiv0$, has no such guarantee. Multiple optimal vertices
and persistent oracle errors can prevent the test from firing.

For a concrete strict improvement, take $Q=[0,1]$ and
$f(t)=t+(t_+)^3$, with $\rho(u)=6u$ on $Q$. At $x=9/20$,
$g=\kappa=643/400$ and $G=5787/8000$. Using only $U=B=G$,
the integrated test passes because $\kappa^2>3G^2$, whereas the
endpoint-profile test fails because $\kappa^2<6G^2$. More generally
their respective thresholds are $x<1/\sqrt3$ and $x<1/\sqrt6$.
These compare sufficient certificates, not one-dimensional
optimization complexity.
